% =====================================================================
%  Foundations of the Lambda-algebra.
%  Sequel to "p-adic Apery limits: rate laws, valuation strata, and closed
%  forms via the p-adic Gamma function" (Paper C of the same series).
%
%  DRAFT.  Author placeholder below.
%  Evidence classes are printed in the statements; see Section 1.4.
%
%  REPAIR RECORD.  This draft incorporates the repairs demanded by the hostile
%  referee pass recorded in REFEREE.md (same directory).  In summary:
%    C-1  the section-3 slogan was stated for c_p modulo the elementary group
%         E_p, which Theorem 3.6 does not give; it is now stated for the tower
%         Betas, with the correct c_p-level corollary and a remark
%         (rem:notclaimed) recording exactly what is NOT claimed.
%    M-1  the "mechanism for Paper C's exclusions" paragraph is now labelled a
%         heuristic and states that Theorem 3.6 does not apply to Lambda_a.
%    M-2/3 the k<->m interchange in Prop 2.3 and the convergence of ghat are now
%         derived from eq (eps), a consequence of the QUOTED convergence of
%         BV eq (2), rather than from an unstated bound on v_p(zeta_p(m)).
%    M-4  u(i,j) is defined in the statement of Thm 4.5, not in its proof.
%    M-5  h_p^{(w)} for general w is no longer attributed to PaperC Thm 6.1
%         (which is w=3); the extension is stated as such.
%    M-6  subsection retitled (it contradicted the remark below it).
%    M-7  Cor 4.10 carries its hypothesis in the statement.
%    M-8  the independence step in Thm 5.4 is now proved (leading symbol).
%    M-9  LLL language: "no vector below the floor was returned", not "excludes".
%    M-10 the planted relation's third entry y is defined.
%    M-11 the cross-family exclusion now carries its range (p=7, 12-21 digits).
%    M-12 Prop 3.1 on Z[1/p] reduced to the integer case.
%    m-1..m-11 as listed in REFEREE.md.
%  Citation pointers into the companion papers were checked against their .aux
%  label tables; PaperC Tab. 6 -> Tab. 16 was the one error found.  The
%  Beukers-Vlasenko Frobenius-structure-constant theorem is cited without an
%  internal number: the fetched arXiv text numbers it 1.4, Paper C cites 1.5.
% =====================================================================
\documentclass[11pt,reqno]{article}

\usepackage[T1]{fontenc}
\usepackage[utf8]{inputenc}
\usepackage[margin=1.15in]{geometry}
\usepackage{amsmath,amssymb,amsthm,mathtools}
\usepackage{booktabs}
\usepackage{array}
\usepackage{enumitem}
\usepackage{xcolor}
\usepackage{microtype}
\usepackage[colorlinks=true,linkcolor=blue!55!black,citecolor=blue!55!black,%
            urlcolor=blue!55!black]{hyperref}

% ---------------------------------------------------------------- macros
\newcommand{\Qp}{\mathbb{Q}_p}
\newcommand{\Zp}{\mathbb{Z}_p}
\newcommand{\QQ}{\mathbb{Q}}
\newcommand{\ZZ}{\mathbb{Z}}
\newcommand{\NN}{\mathbb{N}}
\newcommand{\RR}{\mathbb{R}}
\newcommand{\Gp}{\Gamma_p}
\newcommand{\gtow}{\gamma_p}
\newcommand{\gh}{\widehat{g}_p}
\newcommand{\bet}{\beta_p}
\newcommand{\Atil}{\widetilde{A}}
\newcommand{\Btil}{\widetilde{B}}
\newcommand{\vp}{v_p}
\newcommand{\Pset}{\mathcal{P}}
\newcommand{\Div}{\mathbb{D}}
\newcommand{\El}{\mathcal{E}_p}
\newcommand{\Kalg}{\mathcal{K}_p}
\newcommand{\Lalg}{\mathcal{L}_p}
\DeclareMathOperator{\im}{im}
\DeclareMathOperator{\lcm}{lcm}
\DeclareMathOperator{\cond}{cond}
\DeclareMathOperator{\rank}{rank}
\newcommand{\bin}[2]{\binom{#1}{#2}}
\newcommand{\floor}[1]{\lfloor #1 \rfloor}

% evidence tags
\newcommand{\PROVED}{{\normalfont\textsc{[proved]}}}
\newcommand{\VER}[1]{{\normalfont\textsc{[verified:} #1\textsc{]}}}
\newcommand{\IMPORTED}{{\normalfont\textsc{[imported]}}}
\newcommand{\OPENTAG}{{\normalfont\textsc{[open]}}}

% ---------------------------------------------------------------- theorems
\theoremstyle{plain}
\newtheorem{theorem}{Theorem}[section]
\newtheorem{proposition}[theorem]{Proposition}
\newtheorem{lemma}[theorem]{Lemma}
\newtheorem{corollary}[theorem]{Corollary}
\newtheorem{conjecture}[theorem]{Conjecture}

\theoremstyle{definition}
\newtheorem{finding}[theorem]{Finding}
\newtheorem{definition}[theorem]{Definition}
\newtheorem{hypothesis}[theorem]{Hypothesis}
\newtheorem{example}[theorem]{Example}
\newtheorem{problem}[theorem]{Problem}

\theoremstyle{remark}
\newtheorem{remark}[theorem]{Remark}
\newtheorem{question}[theorem]{Question}

\numberwithin{equation}{section}

% ---------------------------------------------------------------- title
\title{\bfseries Foundations of the $\Lambda$-algebra:\\
the $\Gp$-tower, and the relations among\\ $p$-adic Ap\'ery limits}

% ---------------------------------------------------------------------
% AUTHOR PLACEHOLDER -- to be filled in.
% ---------------------------------------------------------------------
\author{[AUTHOR NAME(S) AND AFFILIATION TO BE SUPPLIED]}
\date{Draft, \today}

\begin{document}
\maketitle

\begin{abstract}
Let $(A,B)$ be a sporadic Ap\'ery-like pair of weight $w$ and character $\chi$.
The $p$-adic Ap\'ery limit
$\Lambda_a=\lim_{s}\chi(p)^sp^{ws}B(ap^s)/A(ap^s)$ exists in $\Qp$, is not a
$\QQ$-affine function of any single $p$-adic zeta value, and is assembled out of
the Kazandzidis limits $c_p(A,B)=\lim_t\bin{Ap^t}{Bp^t}$; this is the content of
the companion Paper~C.  Here we ask what these constants \emph{are} as an
algebra.

We introduce the \emph{$\Gp$-tower product}
$\gtow(x)=\prod_{k\ge1}\Gp(xp^k)$ and the \emph{tower Beta}
$\bet(B,D)=\gtow(B+D)/(\gtow(B)\gtow(D))$, and prove
$c_p(A,B)=E_p(A,B)\,\bet(B,A-B)$, where $E_p$ is an explicit finite monomial in
$p$ and in Morita $\Gp$-values at positive integers.  This single computation
reproves the existence of $c_p(A,B)$, isolates its transcendental part, and
makes the relation problem combinatorial: writing
$\delta(A,B)=[A]-[B]-[A-B]$ in the free abelian group on $\ZZ_{>0}$, we prove
$\log\bet(B,D)$ depends on $(B,D)$ only through $\delta$, that $\delta$ is onto
the augmentation-type subgroup $\ker(\sum c_x\mapsto\sum c_xx)$, and hence
(\textbf{Rigidity Theorem}) that a multiplicative relation among the tower Betas
either lies in the lattice $\ker\delta$ --- generated, in the verified range, by
binomial symmetry and trinomial revision --- or produces a nontrivial vanishing
$p$-adic linear form
$\sum_{m\ \mathrm{odd}\ge3}v_m\zeta_p(m)p^m/(m(1-p^m))=0$ with explicit integer
coefficients.  Under the resulting hypothesis $(\mathrm{IND}_p)$, the group
generated by the tower Betas is free abelian of countably infinite rank.  We
report the exclusion ledger for $(\mathrm{IND}_p)$: nothing below the noise
floor at $p=5,7,11,13$ for $x\le8$ at $26$ and at $38$ certified $p$-adic
digits, with the best heights scaling with the precision, as noise does.

We then define the $\Lambda$-algebra $\Lalg$ --- the closure of the algebra
generated by $\Gp$-values, tower products and the tower-harmonic constants
$h_p(a)$ --- and \emph{prove} the block decomposition of $\Lambda_a$ that
Paper~C established block by block only, thereby proving $\Lambda_a\in\Lalg$ for
the classical Ap\'ery pair and every tower base $a$.  Finally we state the
archimedean/$p$-adic asymmetry precisely: by the Transfer Principle of Paper~D
the two completions share the rational skeleton of the construction exactly,
yet the archimedean Ap\'ery limits span a one-dimensional $\QQ$-space per family
while the $p$-adic ones require, under $(\mathrm{IND}_p)$, infinitely many
independent generators.
\end{abstract}

\tableofcontents

% =====================================================================
\section{Introduction}\label{sec:intro}
% =====================================================================

\subsection{The object}\label{ss:object}

Let $A(n)$ be one of the fifteen sporadic Ap\'ery-like binomial sums
\cite{Zagier2009,Cooper2012,MalikStraub,Straub2023} and $B(n)$ the second
solution of its recurrence, normalised by $B(0)=0$, $B(1)=1$.  Each pair carries
a weight $w\in\{2,3\}$ and a quadratic (possibly trivial) Dirichlet character
$\chi$; when the archimedean Ap\'ery limit $\lim_nB(n)/A(n)$ exists it is a
rational multiple of $\zeta(w)$ or of $L(\chi,w)$
\cite{ASvSZ2008,ChamberlandStraub,PaperA}.  Along a $p$-power tower $n_s=ap^s$
with $p\nmid a$ the normalised ratios converge $p$-adically, and
\begin{equation}\label{eq:Lam}
  \Lambda_a\ :=\ \lim_{s\to\infty}\chi(p)^{s}\,p^{ws}\,\frac{B(ap^s)}{A(ap^s)}\ \in\ \Qp
\end{equation}
is the \emph{$p$-adic Ap\'ery limit}.  Its existence, its convergence rate, the
proof that it is not a $\QQ$-affine function of $\zeta_p(3)$ or $\zeta_p(5)$ at
accessible heights, and a block decomposition of it in terms of the
\emph{Kazandzidis limits}
\begin{equation}\label{eq:cp}
  c_p(A,B)\ :=\ \lim_{t\to\infty}\bin{Ap^t}{Bp^t}\ \in\ \Qp
  \qquad(A>B\ge1)
\end{equation}
are the content of the companion Paper~C \cite{PaperC}.  The present paper takes
those results as given and asks the next question: \emph{what algebra do these
constants generate, and what are its relations?}

\subsection{What this paper adds, and what it imports}\label{ss:position}

This is deliberately a foundations paper for a single class of constants, and
its debts should be stated before its results.

\paragraph{Imported, and used as a black box.}
\begin{itemize}[leftmargin=1.4em]
\item From Paper~A \cite{PaperA}: the twisted descent law $(\mathrm{LB}^\chi_w)$,
  $\vp\bigl(p^wB_nA_q-\chi(p)B_qA_n\bigr)\ge w$ with $q=\floor{n/p}$, which makes
  \eqref{eq:Lam} converge; and the table of the fifteen families with their
  weights, characters and archimedean limits.
\item From Paper~C \cite{PaperC}: the existence theorem for $\Lambda_a$, the rate
  law $\min(3,r(\chi,w))$, the exclusion theorem (no low-height $\QQ$-affine
  relation to $\zeta_p(3)$ or $\zeta_p(5)$), the exact closed form of the
  tower-harmonic constant $h_p(a)$, and the block decomposition of $\Lambda_a$,
  there established block by block as a Finding.
\item From Paper~D \cite{PaperD}: the Transfer Principle --- the archimedean sum
  and the Volkenborn integral of one rational function are linear forms with
  \emph{the same} rational coefficients.  We use it in \S\ref{sec:mirror} and
  reprove nothing of it.
\end{itemize}

\paragraph{New here.}
\begin{enumerate}[label=(\Alph*),leftmargin=2.1em]
\item \emph{The $\Gp$-tower packaging} (\S\ref{sec:tower}).  We introduce
  $\gtow(x)=\prod_{k\ge1}\Gp(xp^k)$ and prove in one computation
  (Theorem~\ref{thm:closed}) both that $c_p(A,B)$ exists and that
  \[
    c_p(A,B)=E_p(A,B)\cdot\bet(B,A-B),\qquad
    \bet(B,D)=\frac{\gtow(B+D)}{\gtow(B)\gtow(D)},
  \]
  with $E_p$ elementary.  This separates the arithmetic of $c_p$ into a finite
  $\Gp$-part and a single transcendental factor, and it is the form on which
  everything downstream runs.  It also discharges, for these constants, an
  orientation remark left open in \cite[\S2.4]{PaperB}.
\item \emph{Rigidity} (\S\ref{sec:relations}).  Multiplicative relations among
  the $\bet$ are governed by a divisor map into a free abelian group
  (Theorem~\ref{thm:rigid}).  Either a relation is \emph{binomial} --- an element
  of the lattice $\ker\delta$, which in the verified range is generated by
  symmetry and trinomial revision --- or it certifies a nontrivial vanishing
  $p$-adic linear form in the odd Kubota--Leopoldt zeta values.
\item \emph{The $\Lambda$-algebra} (\S\ref{sec:algebra}).  We define $\Lalg$ and
  \emph{prove} the block decomposition of $\Atil_a=\lim_sA(ap^s)$ and of
  $\Lambda_a$ for the classical Ap\'ery pair and every $a$ with $p\nmid a$
  (Theorems~\ref{thm:Atil}, \ref{thm:Lam}), upgrading Paper~C's Finding to a
  theorem and giving $\Lambda_a\in\Lalg$.
\item \emph{The asymmetry} (\S\ref{sec:mirror}), stated as a
  precise contrast rather than a slogan: shared rational skeleton,
  one-dimensional archimedean constant space, infinite-rank $p$-adic one.
\item \emph{An exclusion ledger for $(\mathrm{IND}_p)$}
  (\S\ref{sec:ledger}) and a programme (\S\ref{sec:programme}).
\end{enumerate}

\subsection{Two remarks on scope}

The rigidity theorem is about the constants $c_p(A,B)$, which are the building
blocks; it is not a determination of the relations among the $\Lambda_a$
themselves, which are infinite assemblies of blocks.  Paper~C reports that no
$\QQ$-linear relation among $\Lambda_a$'s, across bases or across families, is
visible at accessible heights; the present paper explains why one should not
expect the question to be finite-dimensional, and isolates the one hypothesis
$(\mathrm{IND}_p)$ on which a clean answer for the blocks depends.

Second: $(\mathrm{IND}_p)$ is a $p$-adic independence statement about odd zeta
values.  It is out of reach --- no $\zeta_p(2k+1)$ is known to be irrational for
a single prime $p\ge5$ \cite{Calegari2005,LSZ2025}.  We therefore state it as a
hypothesis, record what it buys, and report the numerical exclusions honestly.

\subsection{Evidence classes}\label{ss:evidence}

Every statement carries its class, printed with it.

\begin{description}[leftmargin=2.2em,style=nextline]
\item[\PROVED] A complete proof is given here.
\item[\IMPORTED] Proved in a cited companion paper or in the literature, and
  used here without reproof.  The citation is given with the statement.
\item[\VER{range}] A finite exact computation over the stated range at the
  stated $p$-adic precision.  This is evidence, never proof.
\item[\OPENTAG] Open, with its sharpest known formulation.
\end{description}

Statements labelled \textsc{theorem}, \textsc{proposition}, \textsc{lemma},
\textsc{corollary} are \PROVED\ or \IMPORTED\ as marked; statements labelled
\textsc{finding} rest on \VER\ evidence only and are never used as inputs to a
proof.  All computations reported here use exact integer or rational arithmetic;
no floating point occurs at any stage, and every $p$-adic digit is derived from
an exact rational or an exact residue.  Programs are listed in
Appendix~\ref{app:repro}.

\subsection{Notation}\label{ss:notation}

$p$ is a prime, always $\ge5$.  $\omega$ is the Teichm\"uller character modulo
$p$ and $\langle x\rangle=x/\omega(x)$; following \cite[Def.~8]{LSZ2025} we set
$\zeta_p(s)=L_p(s,\omega^{1-s})$ for integers $s\ge2$, the same normalisation as
\cite{BV2025}.  In particular $\zeta_p(2k)=0$ for $k\ge1$, used constantly.
$\Gp$ is Morita's $p$-adic Gamma function \cite{Morita1975}, so that for a
positive integer $n$
\begin{equation}\label{eq:morita-int}
  \Gp(n)=(-1)^{n}\!\!\prod_{\substack{1\le j<n\\ p\nmid j}}\!\! j,
  \qquad\text{and}\qquad
  n!=p^{\vp(n!)}\prod_{i\ge0}(-1)^{n_i+1}\Gp(n_i+1),\quad n_i=\floor{n/p^i},
\end{equation}
the product being the iterate of Morita's relation
$\Gp(n+1)=(-1)^{n+1}n!/(p^{\floor{n/p}}\floor{n/p}!)$ and finite because
$\Gp(1)=-1$ makes the terms with $n_i=0$ equal to $1$.  We use throughout the
Beukers--Vlasenko expansion \cite[eq.~(2)]{BV2025}
\begin{equation}\label{eq:BV2}
  \log\Gp(y)\ =\ \Gp'(0)\,y-\sum_{m\ge2}\zeta_p(m)\frac{y^m}{m},
  \qquad y\in p\Zp .
\end{equation}
$\Div:=\ZZ[\ZZ_{>0}]$ is the free abelian group on the symbols $[x]$, $x\ge1$,
and $\Pset:=\{(A,B)\in\ZZ^2:\ A>B\ge1\}$.

\emph{A notational warning, inherited from Paper~C.}  The letters $A,B$ do
double duty: $(A(n),B(n))$ is the sporadic pair of \emph{sequences} of
\S\ref{ss:object}, while in $c_p(A,B)$, $E_p(A,B)$ and $\delta(A,B)$ the
arguments $A,B$ are \emph{positive integers}.  The two never occur in the same
formula, and the sequences are always written with their argument, $A(n)$,
$B(n)$.  We also set $c_p(A,A):=1$, consistent
with $\bin{Ap^t}{Ap^t}=1$; the pairs $(A,A)$ are excluded from $\Pset$ because
they carry no information.

% =====================================================================
\section{The $\Gp$-tower}\label{sec:tower}
% =====================================================================

\subsection{The tower product and the tower Beta}

\begin{lemma}\label{lem:lip}
\IMPORTED\ (i)--(ii), \PROVED\ (iii)--(iv).  Let $p$ be odd.
\begin{enumerate}[label=(\roman*),leftmargin=2em]
\item $\Gp(0)=1$ and $|\Gp(x)-\Gp(y)|_p\le|x-y|_p$ for all $x,y\in\Zp$
  \textup{(}\cite{Morita1975,Robert2000}\textup{)}.
\item $\Gp(1+y)=-\Gp(y)$ for $y\in p\Zp$ \textup{(}same sources\textup{)}.
\item Hence $\Gp(y)\in1+p^{e}\Zp$ and $\vp\bigl(\log\Gp(y)\bigr)\ge e$ whenever
  $\vp(y)=e\ge1$.
\item Hence $\Gp'(0)\in\Zp$.
\end{enumerate}
\end{lemma}

\begin{proof}
(iii): by (i) with $x=0$, $|\Gp(y)-1|_p\le|y|_p=p^{-e}$; the logarithm of an
element of $1+p^e\Zp$ lies in $p^e\Zp$ for $p$ odd.  (iv): by \eqref{eq:BV2},
$\Gp'(0)=\lim_{y\to0}\log\Gp(y)/y$ along $y\in p\Zp$, and by (iii) the quotient
lies in $\Zp$ throughout.
\end{proof}

\begin{definition}\label{def:tower}
For $x\in\Zp$ put
\[
  \gtow(x):=\prod_{k\ge1}\Gp(xp^k)\ \in\ 1+p\Zp,
  \qquad
  \gh(x):=-\!\!\sum_{\substack{m\ge3\\ m\ \mathrm{odd}}}\!\!
    \zeta_p(m)\,\frac{x^m\,p^m}{m\,(1-p^m)} .
\]
We call $\gtow$ the \emph{$\Gp$-tower product} and $\gh$ the \emph{$\Gp$-tower
function}.  For $B,D\ge1$ the \emph{tower Beta} is
\[
  \bet(B,D):=\frac{\gtow(B+D)}{\gtow(B)\,\gtow(D)} .
\]
\end{definition}

The product converges by Lemma~\ref{lem:lip}(iii).  For $\gh$ we use the
following consequence of \eqref{eq:BV2}, and nothing else about $\zeta_p$:
since the series in \eqref{eq:BV2} converges for \emph{every} $y\in p\Zp$,
its coefficients satisfy
\begin{equation}\label{eq:eps}
  \varepsilon_m:=\vp\bigl(\zeta_p(m)/m\bigr)\quad\text{obeys}\quad
  \varepsilon_m+m\longrightarrow\infty .
\end{equation}
The $m$-th term of $\gh(x)$ then has valuation $\ge\varepsilon_m+m$ for
$x\in\Zp$, because $\vp\bigl(p^m/(1-p^m)\bigr)=m$; so $\gh$ converges on $\Zp$.

\begin{proposition}\label{prop:gtow}
\PROVED.  Write $\ell_p:=\Gp'(0)\,p/(1-p)$.  Then for $x\in\Zp$
\begin{equation}\label{eq:loggtow}
  \log\gtow(x)\ =\ \ell_p\,x+\gh(x),
\end{equation}
and consequently, for all $B,D\ge1$,
\begin{equation}\label{eq:logbeta}
  \log\bet(B,D)\ =\ \gh(B+D)-\gh(B)-\gh(D).
\end{equation}
Moreover $\ell_p\in p\Zp$ and $\gh(x)\in p\Zp$ for every $x\in\Zp$, and $\gtow$
satisfies the \emph{distribution relation}
\begin{equation}\label{eq:distrib}
  \gtow(px)=\frac{\gtow(x)}{\Gp(px)},
  \qquad\text{equivalently}\qquad
  \gh(px)=\gh(x)-\log\bigl(\Gp(px)e^{-\Gp'(0)px}\bigr).
\end{equation}
\end{proposition}

\begin{proof}
Substituting $y=xp^k$ in \eqref{eq:BV2} and summing over $k\ge1$ requires
interchanging the two summations.  This is legitimate: by \eqref{eq:eps} the
$(k,m)$ term has valuation
$\ge\varepsilon_m+mk\ge(\varepsilon_m+m)+2(k-1)$, which tends to $\infty$
jointly in $(k,m)$, so the double family is summable and may be summed in either
order.  Using the $p$-adically convergent geometric series
$\sum_{k\ge1}p^{k}=p/(1-p)$ and $\sum_{k\ge1}p^{km}=p^m/(1-p^m)$ this gives
$\log\gtow(x)=\Gp'(0)x\,p/(1-p)-\sum_{m\ge2}\zeta_p(m)x^mp^m/(m(1-p^m))$;
the even $m$ drop out because $\zeta_p(2k)=0$.  This is \eqref{eq:loggtow}.
By Lemma~\ref{lem:lip}(iii) each factor $\Gp(xp^k)$ lies in $1+p^k\Zp$, so
$\gtow(x)\in1+p\Zp$ and $\log\gtow(x)\in p\Zp$; and $\ell_p\in p\Zp$ by
Lemma~\ref{lem:lip}(iv), whence $\gh(x)=\log\gtow(x)-\ell_px\in p\Zp$.
Since $(B+D)-B-D=0$, the linear terms cancel in $\log\bet$, giving
\eqref{eq:logbeta}.  Finally $\gtow(px)=\prod_{k\ge1}\Gp(xp^{k+1})
=\prod_{k\ge2}\Gp(xp^k)=\gtow(x)/\Gp(px)$, and the second form follows on taking
logarithms in \eqref{eq:loggtow}.
\end{proof}

Equation~\eqref{eq:logbeta} is the point of the packaging: \emph{the
transcendental content of the tower depends on the pair only through the formal
difference $[B+D]-[B]-[D]$}.  This is what \S\ref{sec:relations} exploits.

\begin{remark}[relation to Beukers--Vlasenko]\label{rem:BV}
Write $G(y)=\Gp(y)e^{-\Gp'(0)y}=\exp\bigl(-\sum_{m\ge2}\zeta_p(m)y^m/m\bigr)$,
the generating function whose Taylor coefficients at $0$ are the Frobenius
structure constants $\alpha_j$ of \cite{BV2025}.  Then
$e^{\gh(x)}=\prod_{k\ge1}G(xp^k)$: the constants of this paper are
\emph{products of values} of $G$ along a $p$-power tower, where the $\alpha_j$
are \emph{Taylor coefficients} of $G$ at $0$.  For the order-$3$ Ap\'ery
operator the available slots $\alpha_1,\alpha_2$ both vanish
\cite[\S5.4]{PaperC}, so no $\alpha_j$ can be one of our constants.  This
relationship was recorded as an orientation remark, explicitly disclaimed, in
\cite[\S2.4]{PaperB}; Theorem~\ref{thm:closed} discharges it for the
Kazandzidis limits.
\end{remark}

\subsection{The closed form}

For $A>B\ge1$ with $D=A-B$ set $A_l=\floor{A/p^l}$, $B_l=\floor{B/p^l}$,
$D_l=\floor{D/p^l}$ and
\begin{equation}\label{eq:Efin}
  T(A,B):=\prod_{l\ge1}
    \frac{(-1)^{A_l+1}\Gp(A_l+1)}
         {(-1)^{B_l+1}\Gp(B_l+1)\cdot(-1)^{D_l+1}\Gp(D_l+1)},
  \qquad
  E_p(A,B):=-\,p^{\vp\binom{A}{B}}\,T(A,B)\,
    \frac{\Gp(A+1)}{\Gp(B+1)\Gp(D+1)} .
\end{equation}
Both products are finite; $T(A,B)=1$ when $A<p$.  We call $E_p(A,B)$ the
\emph{elementary part}.

\begin{theorem}\label{thm:closed}
\PROVED.  For every prime $p\ge5$ and all integers $A>B\ge1$ the limit
$c_p(A,B)=\lim_t\bin{Ap^t}{Bp^t}$ exists in $\Qp$, and
\begin{equation}\label{eq:closed}
  \boxed{\ c_p(A,B)\ =\ E_p(A,B)\cdot\bet(B,A-B).\ }
\end{equation}
Moreover $\vp\bigl(c_p(A,B)\bigr)=\vp\bin{A}{B}$, and the convergence is at
$3$ digits per level.
\end{theorem}

\begin{proof}
Fix $t\ge1$ and apply \eqref{eq:morita-int} to $Ap^t$, $Bp^t$, $Dp^t$.  By
Kummer's theorem $\vp\bin{Ap^t}{Bp^t}$ counts the carries in adding $Bp^t$ and
$Dp^t$ in base $p$, which is the carry count of $B+D$; hence
$\vp\bin{Ap^t}{Bp^t}=\vp\bin{A}{B}$ for all $t$.  The unit parts give
\[
  \bin{Ap^t}{Bp^t}
  = p^{\vp\binom{A}{B}}\prod_{i\ge0}
     \frac{(-1)^{X_i+1}\Gp(X_i+1)}
          {(-1)^{Y_i+1}\Gp(Y_i+1)\,(-1)^{Z_i+1}\Gp(Z_i+1)},
\]
with $X_i=\floor{Ap^t/p^i}$ and similarly for $Y_i,Z_i$.  Split the product at
$i=t$.

For $i>t$ one has $X_i=A_{i-t}$, $Y_i=B_{i-t}$, $Z_i=D_{i-t}$, so the tail is
exactly $T(A,B)$, independent of $t$.

For $0\le i\le t$ put $k=t-i$, so $X_i=Ap^{k}$, $Y_i=Bp^k$, $Z_i=Dp^k$.  The
signs contribute $(-1)^{Ap^k+1-(Bp^k+1)-(Dp^k+1)}=(-1)^{-1}=-1$ for each $k$,
since $A=B+D$.  For $k\ge1$ the arguments lie in $p\Zp$, where
$\Gp(1+y)=-\Gp(y)$, so
\[
  \frac{\Gp(Ap^k+1)}{\Gp(Bp^k+1)\Gp(Dp^k+1)}
  =\frac{-\Gp(Ap^k)}{(-\Gp(Bp^k))(-\Gp(Dp^k))}
  =-\frac{\Gp(Ap^k)}{\Gp(Bp^k)\Gp(Dp^k)},
\]
and the two minus signs cancel.  The factor $k=0$ contributes
$-\Gp(A+1)/(\Gp(B+1)\Gp(D+1))$.  Altogether
\begin{equation}\label{eq:finite-t}
  \bin{Ap^t}{Bp^t}
  = E_p(A,B)\prod_{k=1}^{t}\frac{\Gp(Ap^k)}{\Gp(Bp^k)\Gp(Dp^k)} .
\end{equation}
By Lemma~\ref{lem:lip} each factor lies in $1+p^k\Zp$ and differs from $1$ by at
most $p^{-k}$, so the partial products form a Cauchy sequence converging to
$\gtow(A)/(\gtow(B)\gtow(D))=\bet(B,D)$.  Hence the limit \eqref{eq:cp} exists
and equals \eqref{eq:closed}.  Since $E_p(A,B)$ is a $p$-power times a unit and
$\bet$ is a unit, $\vp(c_p)=\vp\bin{A}{B}$.

For the rate we quote the classical Jacobsthal--Kazandzidis congruence
$\bin{np}{mp}\equiv\bin nm \pmod{p^{\vp(p^3nm(n-m))}}$
\cite{JacobsthalKazandzidis,Mestrovic}: at $n=Ap^t$, $m=Bp^t$ it gives
\[
  \bin{Ap^{t+1}}{Bp^{t+1}}\ \equiv\ \bin{Ap^{t}}{Bp^{t}}
  \pmod{p^{\,3+3t+\vp(ABD)}},
\]
and since $\vp\bin{Ap^t}{Bp^t}$ is constant in $t$ the relative gain is at least
$3$ digits per level.
\end{proof}

\begin{corollary}\label{cor:cp21}
\PROVED.  $\displaystyle c_p(2,1)=\lim_s\bin{2p^s}{p^s}
=2\exp\Bigl(-\!\!\sum_{m\ge3\ \mathrm{odd}}\!\!\zeta_p(m)\frac{(2^m-2)p^m}{m(1-p^m)}\Bigr)
=2\,\bet(1,1)$.
\end{corollary}
\begin{proof}
$T(2,1)=1$, $\vp\bin21=0$ and $\Gp(3)/\Gp(2)^2=-2$, so $E_p(2,1)=2$;
$\gh(2)-2\gh(1)$ is the displayed series.
\end{proof}

This recovers \cite[Cor.~6.4]{PaperC}.  The new content of
Theorem~\ref{thm:closed} is not the formula but its shape: existence and closed
form come from one computation, and the transcendental factor is a
\emph{single} tower Beta.

\begin{finding}\label{find:closedver}
\VER{$11$ pairs $(A,B)$ $\times$ $p\in\{5,7,11,13\}$, all $12$ certified digits,
zero mismatches; at $p=5$ also at $15$ digits}  Formula~\eqref{eq:closed} was
checked against $c_p(A,B)$ computed directly as a limit of exact binomial
residues, for
$(A,B)=(2,1),(3,1),(3,2),(4,1),(5,2),(7,3),(p,1),(p,2),(p+1,1),(p+1,p),(2p,p)$
--- including cases with carries, with $p\mid A$, and with $\vp(c_p)=1$.  The
two sides are computed by disjoint code paths: the left from binomial residues
via Legendre/Morita splitting, the right from the Kubota--Leopoldt values
$\zeta_p(m)$.
\end{finding}

\subsection{Extension to $\ZZ[1/p]$}

\begin{proposition}\label{prop:homog}
\PROVED.  $c_p(pA,pB)=c_p(A,B)$ for all $A>B\ge1$.  Consequently $c_p$ extends
uniquely to pairs $A>B>0$ in $\ZZ[1/p]$ by $c_p(A,B):=c_p(p^jA,p^jB)$ for any
$j$ with $p^jA,p^jB\in\ZZ$.
\end{proposition}

\begin{proof}
Directly from \eqref{eq:cp}: $c_p(pA,pB)=\lim_t\bin{Ap^{t+1}}{Bp^{t+1}}$, a
subsequence of the defining sequence for $c_p(A,B)$.
\end{proof}

Proposition~\ref{prop:homog} is trivial from the definition but is a nontrivial
identity for the closed form, and it is worth recording that
\eqref{eq:closed} satisfies it on the nose.

\begin{proposition}\label{prop:homog-check}
\PROVED.  Formula~\eqref{eq:closed} is compatible with
Proposition~\ref{prop:homog}; the verification is the distribution relation
\eqref{eq:distrib}.
\end{proposition}

\begin{proof}
$\vp\bin{pA}{pB}=\vp\bin AB$ (Kummer).  Since $\floor{pA/p^l}=A_{l-1}$,
\[
  T(pA,pB)=\frac{(-1)^{A+1}\Gp(A+1)}
                {(-1)^{B+1}\Gp(B+1)(-1)^{D+1}\Gp(D+1)}\,T(A,B)
          =-\,T(A,B)\,\frac{\Gp(A+1)}{\Gp(B+1)\Gp(D+1)},
\]
using $A=B+D$ for the sign, while
$\Gp(pA+1)/(\Gp(pB+1)\Gp(pD+1))=-\Gp(pA)/(\Gp(pB)\Gp(pD))$ by
$\Gp(1+y)=-\Gp(y)$.  By \eqref{eq:distrib},
$\bet(pB,pD)=\bet(B,D)\,\Gp(pB)\Gp(pD)/\Gp(pA)$.  Multiplying the three
displays, all $\Gp$-factors at $p$-multiples and both extra signs cancel and one
is left with $E_p(A,B)\bet(B,D)$.
\end{proof}

\begin{remark}
The tower function itself does \emph{not} descend: $\gh(px)\ne\gh(x)$, and
\eqref{eq:distrib} measures the discrepancy by a single $\Gp$-value.  It is the
combination $E_p\cdot\bet$ that is $p$-homogeneous.  This is the first sign that
the natural invariants of the family are the tower Betas modulo $\Gp$-values,
which is exactly the quotient in which \S\ref{sec:relations} works.
\end{remark}

% =====================================================================
\section{The relations}\label{sec:relations}
% =====================================================================

\subsection{Three binomial relations}

\begin{proposition}\label{prop:binrels}
\PROVED.  For $c_p$ extended as in Proposition~\ref{prop:homog}:
\begin{enumerate}[label=(\roman*),leftmargin=2em]
\item (boundary) $c_p(A,A)=1$;
\item (symmetry) $c_p(A,B)=c_p(A,A-B)$;
\item ($p$-homogeneity) $c_p(pA,pB)=c_p(A,B)$;
\item (trinomial revision) $c_p(A,B)\,c_p(B,C)=c_p(A,C)\,c_p(A-C,B-C)$ for
  $A>B>C\ge1$.
\end{enumerate}
\end{proposition}

\begin{proof}
Each is a binomial identity that holds \emph{at every level of the tower
simultaneously}, hence passes to the limit: $\bin{N}{K}=\bin{N}{N-K}$ for (ii),
and
$\bin{N}{K}\bin{K}{L}=\bin{N}{L}\bin{N-L}{K-L}$ at $N=Ap^t$, $K=Bp^t$,
$L=Cp^t$ for (iv).  (iii) is Proposition~\ref{prop:homog}, (i) is trivial.  For
arguments in $\ZZ[1/p]$, multiply all of $A,B,C$ by a common $p^j$ clearing
denominators and apply (iii).
\end{proof}

\begin{finding}\label{find:rels}
\VER{$p\in\{5,7,11,13\}$; $9$ certified digits; $66$ symmetry cells
($2\le A\le12$), $21$ $p$-homogeneity cells ($2\le A\le7$) and $120$
trinomial-revision cells ($3\le A\le10$) per prime; $828$ checks in total, zero
failures}  Proposition~\ref{prop:binrels} holds on the \emph{directly computed}
limits, i.e.\ without going through \eqref{eq:closed}.
\end{finding}

That these are verified as well as proved is not redundancy: the direct limits
and the closed form are computed by disjoint code, and the relations are the
only place where the two are compared across many pairs at once.

\subsection{The divisor map}

\begin{definition}\label{def:delta}
Define the group homomorphisms
\[
  \delta:\ZZ[\Pset]\longrightarrow\Div,\quad
  (A,B)\longmapsto[A]-[B]-[A-B];
  \qquad
  m_r:\Div\longrightarrow\ZZ,\quad \textstyle\sum_xc_x[x]\longmapsto\sum_xc_xx^r
  \ (r\ge1);
\]
\[
  \widehat{G}:\Div\longrightarrow\Qp,\quad
  \textstyle\sum_xc_x[x]\longmapsto\sum_xc_x\,\gh(x).
\]
\end{definition}

By \eqref{eq:logbeta}, for any $\mathbf n=\sum_in_i(A_i,B_i)\in\ZZ[\Pset]$,
\begin{equation}\label{eq:key}
  \prod_i\bet(B_i,A_i-B_i)^{n_i}\ =\ \exp\bigl(\widehat G(\delta(\mathbf n))\bigr).
\end{equation}
Everything in this section is a consequence of \eqref{eq:key}.

\begin{proposition}\label{prop:image}
\PROVED.  $\im\delta=\ker m_1$, a subgroup of $\Div$ of infinite rank and
corank~$1$.
\end{proposition}

\begin{proof}
$m_1\delta(A,B)=A-B-(A-B)=0$, so $\im\delta\subseteq\ker m_1$.  Conversely
$\delta(2,1)=[2]-2[1]$, and $\delta(n,1)=[n]-[1]-[n-1]$ gives inductively
$[n]-n[1]=\delta(n,1)+\bigl([n-1]-(n-1)[1]\bigr)\in\im\delta$ for all $n\ge2$.
Finally any $\mu=\sum_xc_x[x]$ with $m_1(\mu)=0$ equals
$\sum_xc_x([x]-x[1])$, hence lies in $\im\delta$.
\end{proof}

\begin{lemma}\label{lem:moments}
\PROVED.  The map $\mu\mapsto(m_3(\mu),m_5(\mu),m_7(\mu),\dots)$ is injective on
$\Div$.  Indeed if $m_r(\mu)=0$ for infinitely many $r$ then $\mu=0$.
\end{lemma}

\begin{proof}
Suppose $\mu=\sum_xc_x[x]\ne0$ and let $X$ be the largest $x$ with $c_x\ne0$.
Then $m_r(\mu)/X^r=c_X+\sum_{x<X}c_x(x/X)^r\to c_X\ne0$ as $r\to\infty$ in $\RR$,
so $m_r(\mu)\ne0$ for all large $r$.
\end{proof}

\subsection{The rigidity theorem}

\begin{theorem}[Rigidity]\label{thm:rigid}
\PROVED.  Let $\mathbf n=\sum_{i=1}^{r}n_i(A_i,B_i)\in\ZZ[\Pset]$ and put
$\mu=\delta(\mathbf n)\in\Div$, $v_m=m_m(\mu)\in\ZZ$.
\begin{enumerate}[label=(\roman*),leftmargin=2em]
\item If $\mu=0$ then, for \emph{every} prime $p\ge5$,
  \[
    \prod_i\bet(B_i,A_i-B_i)^{n_i}=1,
    \qquad\text{hence}\qquad
    \prod_ic_p(A_i,B_i)^{n_i}=\prod_iE_p(A_i,B_i)^{n_i},
  \]
  an explicit rational number times a monomial in $\Gp$-values at positive
  integers.
\item If $\mu\ne0$ and $\prod_i\bet(B_i,A_i-B_i)^{n_i}=1$ for some $p\ge5$, then
  \begin{equation}\label{eq:forced}
    \sum_{\substack{m\ge3\\ m\ \mathrm{odd}}}
      v_m\,\zeta_p(m)\,\frac{p^m}{m(1-p^m)}\ =\ 0
  \end{equation}
  is a nontrivial vanishing $p$-adic linear form in the odd Kubota--Leopoldt
  zeta values: the integers $v_m$ are not all zero, and
  $|v_m|\le\bigl(\sum_x|c_x|\bigr)X^m$ where $X$ is the largest element of the
  support of $\mu$.
\end{enumerate}
\end{theorem}

\begin{proof}
(i) is \eqref{eq:key} with $\widehat G(0)=0$, together with
Theorem~\ref{thm:closed}.  For (ii): by Proposition~\ref{prop:gtow},
$\widehat G(\mu)\in p\Zp$, and $\exp$ is injective on $p\Zp$ for $p$ odd, so
\eqref{eq:key} gives $\widehat G(\mu)=0$.  Each $\gh(x)$ is by
Definition~\ref{def:tower} a convergent series in $\Qp$, and $\mu$ has finite
support, so the finite sum $\sum_xc_x\gh(x)$ may be resummed term by term:
\[
  \widehat G(\mu)
  =-\!\!\sum_{\substack{m\ge3\\ m\ \mathrm{odd}}}\!\!
     \Bigl(\sum_xc_xx^m\Bigr)\zeta_p(m)\frac{p^m}{m(1-p^m)}
  =-\!\!\sum_{\substack{m\ge3\\ m\ \mathrm{odd}}}\!\!
     v_m\,\zeta_p(m)\frac{p^m}{m(1-p^m)} .
\]
Lemma~\ref{lem:moments} says the $v_m$ are not all zero, and
$|v_m|=|\sum_xc_xx^m|\le(\sum_x|c_x|)X^m$.
\end{proof}

\begin{hypothesis}[$\mathrm{IND}_p$]\label{hyp:ind}
\OPENTAG.  The only $\mu\in\Div$ with $\widehat G(\mu)=0$ is $\mu=0$;
equivalently, the values $\gh(1),\gh(2),\gh(3),\dots$ are $\QQ$-linearly
independent; equivalently, \eqref{eq:forced} has no solution in integers $v_m$
of the form $v_m=m_m(\mu)$, $\mu\ne0$.
\end{hypothesis}

\begin{corollary}\label{cor:free}
\PROVED, under Hypothesis~\ref{hyp:ind}.  The multiplicative group generated by
$\{\bet(B,D):B,D\ge1\}$ inside $1+p\Zp$ is free abelian of countably infinite
rank, and $\delta$ induces an isomorphism onto $\ker m_1\subset\Div$.
Unconditionally it is a quotient of $\ker m_1$.
\end{corollary}

\begin{proof}
By \eqref{eq:key} the group is the image of $\ZZ[\Pset]$ under
$\exp\circ\widehat G\circ\delta$.  Its relation module contains $\ker\delta$
always, and equals $\ker\delta$ exactly when $\ker\widehat G\cap\im\delta=0$,
which is Hypothesis~\ref{hyp:ind}.  Then the group is
$\ZZ[\Pset]/\ker\delta\cong\im\delta=\ker m_1$ by
Proposition~\ref{prop:image}.  Finally $\ker m_1$ is free, being a subgroup of
the free abelian group $\Div$, and of infinite rank, since it contains the
independent family $\{[n]-n[1]\}_{n\ge2}$.
\end{proof}

\subsection{The lattice of binomial relations}

Theorem~\ref{thm:rigid}(i) makes $\ker\delta$ the lattice of \emph{expected}
relations.  It contains two visible families, coming from
Proposition~\ref{prop:binrels}(ii) and (iv):
\begin{align}
  (\mathrm{S}_{A,B})&:\quad (A,B)-(A,A-B), \label{eq:relS}\\
  (\mathrm{T}_{A,B,C})&:\quad (A,B)+(B,C)-(A,C)-(A-C,B-C),\qquad A>B>C\ge1.
  \label{eq:relT}
\end{align}
Both are in $\ker\delta$: for \eqref{eq:relT},
$\delta(A,B)+\delta(B,C)=[A]-[A-B]-[C]-[B-C]=\delta(A,C)+\delta(A-C,B-C)$.
Note that $p$-homogeneity is \emph{not} of this type --- $\delta(pA,pB)\ne
\delta(A,B)$ --- consistently with the fact that it moves $\bet$ by a
$\Gp$-value (Proposition~\ref{prop:homog-check}) rather than fixing it.

\begin{finding}\label{find:lattice}
\VER{$N\le11$; exact integer linear algebra, Smith normal forms over $\ZZ$}
Let $\Pset_N=\{(A,B):N\ge A>B\ge1\}$ and $\delta_N$ the restriction of $\delta$.
Then
\[
  \rank\delta_N=N-1,\qquad
  \rank\ker\delta_N=\binom{N}{2}-(N-1)=\binom{N-1}{2},
\]
and $\ker\delta_N$ is generated \emph{over $\ZZ$} by the relations
\eqref{eq:relS} and \eqref{eq:relT} with all indices $\le N$.  The argument is:
the sublattice $L_N$ they span is contained in $\ker\delta_N$ (checked
symbolically, not assumed); it has the same rank, hence the same $\QQ$-span; and
the Smith normal forms of both have all elementary divisors equal to $1$, so both
are saturated in $\ZZ[\Pset_N]$.  Two saturated sublattices with the same
$\QQ$-span coincide.
\end{finding}

The first assertion is Proposition~\ref{prop:image} restricted, hence
\PROVED; the generation statement is \VER{} only.

\begin{problem}\label{prob:gen}
\OPENTAG.  Prove that $\ker\delta$ is generated by \eqref{eq:relS} and
\eqref{eq:relT}.  Equivalently: give a presentation of the free abelian group
$\ker m_1$ by the symbols $(A,B)$ and those two families of relations.
\end{problem}

\subsection{What the theorem says, in words}

Combining: the tower Betas satisfy the binomial relations \eqref{eq:relS},
\eqref{eq:relT} --- and, under Hypothesis~\ref{hyp:ind}, \emph{nothing else}.
Equivalently, at the level of the Kazandzidis limits: the only relations
\[
  \prod_ic_p(A_i,B_i)^{n_i}\ =\ \prod_iE_p(A_i,B_i)^{n_i}
\]
--- those in which the transcendental factor cancels \emph{identically} --- are
the binomial ones.  A single unexpected multiplicative relation among the tower
Betas would, by Theorem~\ref{thm:rigid}(ii), produce a nontrivial vanishing
$\QQ$-linear form in $\zeta_p(3),\zeta_p(5),\zeta_p(7),\dots$ with explicit
integer coefficients of controlled growth.  Since no odd $p$-adic zeta value is
known to be irrational for any $p\ge5$, that would be a substantial result about
$p$-adic $L$-values.  The relation problem for the building blocks of the
$\Lambda$-algebra therefore sits directly on top of the $p$-adic odd-zeta
problem, and not beside it.

\begin{remark}[what is \emph{not} claimed]\label{rem:notclaimed}
Theorem~\ref{thm:rigid} does not determine when $\prod_ic_p(A_i,B_i)^{n_i}$
merely \emph{lies in} $\El$.  That would require knowing
$\El\cap(1+p\Zp)$, which contains $\QQ^\times\cap(1+p\Zp)$ and is therefore
large; the transcendental factor can in principle be absorbed by an elementary
one without vanishing.  Everything we prove is about the tower Betas, where the
elementary part has been divided out by construction.
\end{remark}

% =====================================================================
\section{The $\Lambda$-algebra}\label{sec:algebra}
% =====================================================================

\subsection{Definition}

\begin{definition}\label{def:algebras}
Fix $p\ge5$.  Let
\begin{align*}
  \El&:=\text{the subgroup of }\Qp^\times\text{ generated by }\QQ^\times
        \text{ and }\{\Gp(n):n\in\ZZ_{\ge1}\}, &&\text{(the \emph{elementary group})}\\
  \Kalg&:=\text{the }\QQ\text{-subalgebra of }\Qp\text{ generated by }
        \El\cup\{\bet(B,D):B,D\ge1\}, &&\text{(the \emph{Kazandzidis algebra})}\\
  \Lalg&:=\text{the closure of }\Kalg\bigl[h_p^{(w)}(a):a\ge1,\ w\ge2\bigr]
        \text{ in }\Qp, &&\text{(the \emph{$\Lambda$-algebra})}
\end{align*}
where
\[
  h_p^{(w)}(a):=\lim_{s\to\infty}p^{ws}\!\!\sum_{k\le ap^s}\!\!k^{-w}
   \ =\!\!\sum_{\substack{x\in\ZZ[1/p]\\ 0<x\le a}}\!\!x^{-w}
\]
is the \emph{tower-harmonic constant}.  For $w=3$ this is
\cite[Thm.~6.1]{PaperC}; the proof there --- split $\sum_{k\le n}k^{-w}$ at
$p\mid k$, iterate, and re-index $x=k/p^i$ --- is verbatim the same for every
$w\ge2$, and we use it in that generality.  We write $h_p=h_p^{(3)}$.
\end{definition}

By Theorem~\ref{thm:closed}, $c_p(A,B)\in\Kalg$ for all $A>B\ge1$.  $\Lalg$ is a
closed $\QQ$-subalgebra of $\Qp$; it contains a multiplicative group of infinite
rank (Proposition~\ref{prop:infinite} under Hypothesis~\ref{hyp:ind}).

It is worth saying which $\Gp$-values are and are not in play.  The $\Gp$-values
generating $\El$ are at \emph{positive integers}, where $\Gp$ is an explicit
product of units by \eqref{eq:morita-int}, and at $p$-power arguments, where it
is the tower product $\gtow$.  The $\Gp$-values at rational arguments with small
denominator, by contrast, are \emph{algebraic}: $\Gp(1/2)^2=\pm1$ by reflection
and $\Gp(1/3),\Gp(1/4)$ are Gauss-sum values by Gross--Koblitz
\cite{GrossKoblitz1979}.  Those contribute nothing new to $\Lalg$, and are not
where the transcendence of these constants can live; the tower is.

\begin{proposition}\label{prop:hp}
\PROVED.  $h_p^{(w)}(pa)=p^{-w}h_p^{(w)}(a)$ for every $a>0$ in $\ZZ[1/p]$, and
$h_p^{(w)}$ extends to $\ZZ[1/p]_{>0}$ by this relation.
\end{proposition}
\begin{proof}
$x\mapsto x/p$ is a bijection from $\{x\in\ZZ[1/p]:0<x\le pa\}$ onto
$\{y\in\ZZ[1/p]:0<y\le a\}$, and $x^{-w}=p^{-w}y^{-w}$.
\end{proof}

This is the tower-harmonic shadow of the scaling law
$\Lambda_{p^ka}=\chi(p)^kp^{-wk}\Lambda_a$ of \cite[Prop.~2.5(iii)]{PaperC}.

\begin{remark}[the index range, which matters]\label{rem:index}
Paper~C's series for $h^{(w)}_p$ reads
$h^{(w)}_p(a)=\sum_ip^{wi}S_i(a)$, $S_i(a)=\sum_{k\le ap^i,\,p\nmid k}k^{-w}$,
with $i$ running over \emph{all} integers with $ap^i\ge1$.  For $1\le a<p$ this
is $i\ge0$ and the two descriptions of $h^{(w)}_p(a)$ agree term by term; for
$a>p$ the terms with $i<0$ are exactly the $x$ in the sum with $\vp(x)>0$, and
omitting them breaks Proposition~\ref{prop:hp}.  We record it because
Proposition~\ref{prop:hp} evaluates $h^{(w)}_p$ outside the range $1\le a<p$ in
which the two descriptions are visibly the same, and because the omission is an
easy one to make: an implementation restricted to $i\ge0$ passes every test at
$a<p$ and fails Proposition~\ref{prop:hp} at $a=p$, which is how
Finding~\ref{find:hp}(a) was used here.
\end{remark}

\begin{finding}\label{find:hp}
\VER{$p\in\{5,7\}$, $w=3$; $a\le5$; at $12$ and again at $18$ certified
$p$-adic digits}  (a) Proposition~\ref{prop:hp} holds numerically at $a=1,2,3$.
(b) No $\QQ$-affine relation among $h_p(1),\dots,h_p(N)$, $N\le5$, was found
below the noise floor; every best height sits at the floor and scales with the
precision --- e.g.\ at $p=5$, $N=1$ the best height moves from
$1.6\times10^{4}$ (floor $1.6\times10^{4}$) to $1.6\times10^{6}$ (floor
$2.0\times10^{6}$) when the precision goes from $12$ to $18$ digits, and at
$p=7$, $N=3$ from $1.7\times10^{2}$ (floor $3.4\times10^{2}$) to
$4.0\times10^{3}$ (floor $6.4\times10^{3}$).  Paper~C excludes $h_p(1)\in\QQ$
and $h_p(1)$ quadratic over $\QQ$ \cite[\S5.2]{PaperC}; the relations
\emph{among} different tower bases are new here.
\end{finding}

\subsection{The structure theorem, proved}

Paper~C establishes the block decomposition of $\Lambda_a$ for $a=1$ block by
block, as a Finding.  Here we prove it, for every tower base.  Throughout this
subsection $(a_n,b_n)$ is the classical Ap\'ery $\zeta(3)$ pair,
$a_n=\sum_k\bin nk^2\bin{n+k}k^2$ and $b_n$ the second solution with $b_0=0$,
$b_1=6$, so $w=3$ and $\chi=1$.

\begin{lemma}\label{lem:kummer}
\PROVED.  Fix $p\nmid a$ and set $\kappa=\kappa(a,p):=\floor{\log_p(2a)}+2$.
\begin{enumerate}[label=(\roman*),leftmargin=2em]
\item For all $N\ge K\ge0$, $\ \vp\bin NK\le\floor{\log_pN}+1$.
\item For $i\ge0$ and $1\le j\le ap^i$ with $p\nmid j$,
  $\ i\le\vp\bin{ap^i}{j}\le i+\kappa-1$ and
  $\vp\bin{ap^i+j}{j}\le i+\kappa-1$.
\item $\vp\bin{ap^s}{jp^{s-i}}=\vp\bin{ap^i}{j}$ and
  $\vp\bin{ap^s+jp^{s-i}}{jp^{s-i}}=\vp\bin{ap^i+j}{j}$ for every $s\ge i$.
\end{enumerate}
\end{lemma}

\begin{proof}
By Kummer's theorem \cite[\S1]{Granville1997} $\vp\bin NK$ is the number of
carries in adding $K$ and $N-K$ in base $p$.  (i): all carries occur in
positions $0,\dots,\floor{\log_pN}$, since $N$ has $\floor{\log_pN}+1$ digits.
(iii): multiplying $N$ and $K$ by $p^{s-i}$ shifts all digits and changes no
carry; note $ap^s+jp^{s-i}=(ap^i+j)p^{s-i}$.  (ii), lower bound: take $N=ap^i$,
$K=j$; the last $i$ base-$p$ digits of $N$ are $0$ while $K$ has a nonzero units
digit, so a carry occurs out of position $0$ and, inductively, out of every
position $<i$.  Upper bounds: apply (i) with $N=ap^i$ and with
$N=ap^i+j\le2ap^i$, both of which are $<2ap^{i}\cdot p$, so
$\floor{\log_pN}+1\le i+\floor{\log_p(2a)}+1=i+\kappa-1$.
\end{proof}

\begin{theorem}\label{thm:Atil}
\PROVED.  Let $p\ge5$ and $p\nmid a$.  The limit $\Atil_a:=\lim_sa_{ap^s}$
exists in $\Zp$ and
\begin{equation}\label{eq:Atilblocks}
  \Atil_a\ =\ 1+\sum_{i\ge0}\ \sum_{\substack{1\le j\le ap^i\\ p\nmid j}}
    \Bigl[c_p\bigl(ap^i,j\bigr)\,c_p\bigl(ap^i+j,\,j\bigr)\Bigr]^{2},
\end{equation}
with the double series convergent in $\Zp$.  We write
\begin{equation}\label{eq:udef}
  u(i,j):=\bigl[c_p(ap^i,j)\,c_p(ap^i+j,j)\bigr]^{2},
  \qquad \vp\bigl(u(i,j)\bigr)\ge2i,
\end{equation}
for the $(i,j)$ block.  (For $i=0$, $j=a$ the first factor is $c_p(a,a)=1$.)
\end{theorem}

\begin{proof}
Fix $s$ and split the sum $a_{ap^s}=\sum_{k=0}^{ap^s}A(ap^s,k)$,
$A(n,k)=\bin nk^2\bin{n+k}k^2$, according to $k=0$ and $k=jp^{s-i}$ with
$p\nmid j$, $i=s-\vp(k)\in[0,s]$ and $1\le j\le ap^{i}$.  This is a bijection
between $\{1,\dots,ap^s\}$ and $\{(i,j):0\le i\le s,\ 1\le j\le ap^i,\
p\nmid j\}$, and $A(ap^s,0)=1$.

Write $u_s(i,j):=A(ap^s,jp^{s-i})$.  By Lemma~\ref{lem:kummer},
$\vp\bin{ap^s}{jp^{s-i}}=\vp\bin{ap^i}{j}\ge i$, so
\begin{equation}\label{eq:unifbd}
  \vp\bigl(u_s(i,j)\bigr)\ \ge\ 2i\qquad\text{for all }s\ge i,
\end{equation}
uniformly.  Moreover, since $ap^s=(ap^i)p^{s-i}$ and
$ap^s+jp^{s-i}=(ap^i+j)p^{s-i}$, Theorem~\ref{thm:closed} gives
\begin{equation}\label{eq:blocklimit}
  \lim_{s\to\infty}u_s(i,j)
  =\bigl[c_p(ap^i,j)\,c_p(ap^i+j,j)\bigr]^{2}=u(i,j),
\end{equation}
and $\vp(u(i,j))\ge2i$ by the same bound, which is \eqref{eq:udef}.

Now fix $I\ge1$.  For every $s\ge I$, by \eqref{eq:unifbd},
\[
  \Bigl|\,a_{ap^s}-1-\sum_{i<I}\sum_ju_s(i,j)\,\Bigr|_p\ \le\ p^{-2I},
  \qquad
  \Bigl|\sum_{i\ge I}\sum_ju(i,j)\Bigr|_p\ \le\ p^{-2I},
\]
the second because a $p$-adic series with terms of valuation $\ge2I$ converges
and has valuation $\ge2I$.  The inner sum $\sum_{i<I}\sum_j$ is finite, so by
\eqref{eq:blocklimit} there is $S(I)$ with
$|\sum_{i<I}\sum_j(u_s(i,j)-u(i,j))|_p\le p^{-2I}$ for $s\ge S(I)$.  Combining
the three estimates, $|a_{ap^s}-\Sigma|_p\le p^{-2I}$ for $s\ge\max(I,S(I))$,
where $\Sigma$ denotes the right-hand side of \eqref{eq:Atilblocks} (which
converges, again by the valuation bound).  As $I$ was arbitrary, $a_{ap^s}\to
\Sigma$.
\end{proof}

\begin{remark}
The proof does not use the Beukers--Coster supercongruence
$a_{mp^r}\equiv a_{mp^{r-1}}\pmod{p^{3r}}$ \cite{Supercong}; existence of
$\Atil_a$ falls out of the block estimate.  It does not recover the $p^{3r}$
either: \eqref{eq:unifbd} gives only $2$ digits per level, and the sharp rate is
a genuinely finer fact.
\end{remark}

\begin{theorem}\label{thm:Lam}
\PROVED.  Let $p\ge5$ and $p\nmid a$.  Then $\Btil_a:=\lim_sp^{3s}b_{ap^s}$
exists and
\begin{equation}\label{eq:Btil}
  \Btil_a=h_p(a)\,\Atil_a
   +\sum_{(i,j)}\sum_{(i',j')}\ (-1)^{j'-1}\,
     \frac{u(i,j)\,p^{3i'}}{2\,j'^{\,3}\,c_p(ap^{i'},j')\,c_p(ap^{i'}+j',j')},
\end{equation}
the inner sum being restricted to the pairs with $j'p^{-i'}\le jp^{-i}$, with
the $((i,j),(i',j'))$ block of valuation $\ge2i+i'-2\kappa+2$, with
$\kappa=\kappa(a,p)$ as in Lemma~\ref{lem:kummer}, and both series convergent.  In particular, if
$\Atil_a\ne0$ then
\begin{equation}\label{eq:Lamblocks}
  \Lambda_a=\frac{\Btil_a}{\Atil_a}=h_p(a)+\frac{1}{\Atil_a}\sum_{(i,j),(i',j')}(\cdots)
  \ \in\ \Lalg .
\end{equation}
\end{theorem}

\begin{proof}
Ap\'ery's formula \cite{Apery1979,Fischler2003} reads
$b_n=\sum_kA(n,k)\bigl[H_3(n)+T(n,k)\bigr]$ with $H_3(n)=\sum_{k\le n}k^{-3}$
and $T(n,k)=\sum_{m\le k}(-1)^{m-1}\bigl(2m^3\bin nm\bin{n+m}m\bigr)^{-1}$.
Hence
\[
  p^{3s}b_{ap^s}=\bigl(p^{3s}H_3(ap^s)\bigr)\,a_{ap^s}
   +\sum_k A(ap^s,k)\sum_{m\le k}\frac{(-1)^{m-1}\,p^{3s}}
     {2m^3\bin{ap^s}{m}\bin{ap^s+m}{m}} .
\]
The first summand converges to $h_p(a)\Atil_a$ by \cite[Thm.~6.1]{PaperC} and
Theorem~\ref{thm:Atil}.  For the second, index $k=jp^{s-i}$ and $m=j'p^{s-i'}$
as before.  Then $\vp(m^3)=3(s-i')$ and, by Lemma~\ref{lem:kummer}(ii)--(iii),
$\vp\bin{ap^s}{m}=\vp\bin{ap^{i'}}{j'}\le i'+\kappa-1$ and likewise
$\vp\bin{ap^s+m}{m}=\vp\bin{ap^{i'}+j'}{j'}\le i'+\kappa-1$, so the inner term
\[
  w_s(i',j'):=\frac{p^{3s}}{2\,m^3\bin{ap^s}{m}\bin{ap^s+m}{m}}
  \qquad\text{has}\qquad \vp\bigl(w_s(i',j')\bigr)\ \ge\ i'-2\kappa+2,
\]
uniformly in $s$, and by Theorem~\ref{thm:closed}
$w_s(i',j')\to p^{3i'}\bigl(2j'^{\,3}c_p(ap^{i'},j')c_p(ap^{i'}+j',j')\bigr)^{-1}$.
Together with \eqref{eq:unifbd} the $((i,j),(i',j'))$ block has valuation
$\ge2i+i'-2\kappa+2$, which tends to $\infty$ jointly in $(i,i')$; the same
three-estimate argument as in Theorem~\ref{thm:Atil}, now with a double
truncation, gives \eqref{eq:Btil}.  Finally every ingredient of
\eqref{eq:Btil} lies in $\Kalg[h_p(a)]$ and the series converges in $\Qp$, so
$\Btil_a$ and, when $\Atil_a\ne0$, $\Lambda_a$ lie in $\Lalg$.
\end{proof}

\begin{remark}[on $\Atil_a\ne0$]\label{rem:Atilnonzero}
$\Lambda_a$ exists unconditionally by \cite[Thm.~2.4]{PaperC}, so
$\Lambda_a\Atil_a=\Btil_a$ holds whenever $\Atil_a\ne0$, and $\Atil_a\ne0$ as
soon as $\vp(\Atil_a)<\infty$.  \VER{$p\in\{5,7\}$, $a\in\{1,2\}$:
$\vp(\Atil_a)\in\{0,1\}$; and $\vp(\Atil_1)\le1$ at $p=5,7,11,13$
\cite[Tab.~16]{PaperC}}, so the hypothesis is harmless in the computed range;
we do not know it in general.  This is the only conditional step in
\eqref{eq:Lamblocks}.
\end{remark}

\begin{corollary}\label{cor:inL}
\PROVED.  Let $p\ge5$, $p\nmid a$, and assume $\Atil_a\ne0$ \textup{(}see
Remark~\textup{\ref{rem:Atilnonzero}}: this is not known in general, and is
\VER{} only on the cells listed there\textup{)}.  Then, for the classical
Ap\'ery pair, $\Lambda_a\in\Lalg$.
\end{corollary}

\begin{finding}\label{find:blocks}
\VER{$p\in\{5,7\}$; $a=1,2$; $12$ certified digits; $\Atil_a$ computed
independently from the exact integer double sum at $n=ap^s$, $s\le4$ resp.\ $3$}
Truncating \eqref{eq:Atilblocks} at $i\le I$ leaves a residual of valuation
$3$--$4$ for $I=0$ and $6$--$7$ for $I=1$, against the guaranteed $2(I+1)$; every
predicted block is confirmed and the residual moves to the next block.  Paper~C
records the case $a=1$; the case $a=2$ is new here and is the numerical
companion to the general-$a$ statement of Theorem~\ref{thm:Atil}.
\end{finding}

\subsection{Why no finite set of tower Betas suffices}

\begin{proposition}\label{prop:infinite}
\PROVED, under Hypothesis~\ref{hyp:ind}.  For every $N\ge2$ the group generated
by $\{\bet(B,D):B+D\le N\}$ is free abelian of rank exactly $N-1$, and the
lattice of relations among the images of the $\binom N2$ symbols of $\Pset_N$
has rank $\binom{N-1}{2}$.  Consequently no finitely generated subgroup of $1+p\Zp$
contains all the tower Betas.
\end{proposition}
\begin{proof}
The generators are indexed by $\Pset_N$, and by \eqref{eq:key} the group is the
image of $\ZZ[\Pset_N]$ under $\exp\circ\widehat G\circ\delta$.  Under
Hypothesis~\ref{hyp:ind} the kernel of that map is $\ker\delta_N$, so the group
is isomorphic to $\im\delta_N$, of rank $N-1$ by
Proposition~\ref{prop:image}; the relation lattice then has rank
$\binom N2-(N-1)=\binom{N-1}2$.  The ranks are unbounded in $N$, so no finitely
generated subgroup contains all the $\bet$.
\end{proof}

\begin{remark}
We do \emph{not} deduce that $\Kalg$ is not a finitely generated $\QQ$-algebra:
a finitely generated algebra can perfectly well contain a multiplicative group
of infinite rank (as $\QQ$ does).  The correct statement is the one about
groups, and it is the one the asymmetry of \S\ref{sec:mirror} uses.
\end{remark}

The decomposition \eqref{eq:Atilblocks} uses $c_p(ap^i,j)$ for all $i$, hence
tower Betas at unboundedly large arguments.  Under
Hypothesis~\ref{hyp:ind} these are independent, so the $p$-adic Ap\'ery limit is
not merely \emph{an} infinite sum: it is an infinite sum over a basis.  This is
the precise sense in which the one-generator hypothesis refuted numerically in
\cite[\S6.6]{PaperC} fails structurally.

% =====================================================================
\section{The mirror}\label{sec:mirror}
% =====================================================================

\subsection{What the two completions share}

\begin{theorem}[Transfer Principle; Paper~D \cite{PaperD}]\label{thm:transfer}
\IMPORTED.  Let $R\in\QQ(t)$ have poles only at non-positive integers and
$\deg R\le-2$, with partial fractions $R(t)=\sum_{i,k}r_{i,k}(t+k)^{-i}$; set
$\rho_i=\sum_kr_{i,k}$ and
$\rho_{0,\theta}=-\sum_{i,k}r_{i,k}\sum_{\nu=0}^{k-1}(\nu+\theta)^{-i}$.  Then
\[
  \sum_{m\ge0}R(m+\theta)=\rho_{0,\theta}+\sum_{i\ge2}\rho_i\,\zeta(i,\theta),
  \qquad
  -\int_{\Zp}\widetilde R(t+\theta)\,dt
   =\rho_{0,\theta}+\sum_{i\ge2}\rho_i\,\omega(\theta)^{1-i}\zeta_p(i,\theta),
\]
with the \emph{same} rational numbers $\rho$, $\widetilde R$ a primitive of $R$.
Consequently every transfer identity and every integrality lemma of
Rhin--Viola type, being a statement about the $\rho$'s, holds verbatim in both
completions.
\end{theorem}

The Transfer Principle says that the \emph{rational skeleton} of an Ap\'ery-type
construction --- which linear form in which zeta values, with which coefficients
--- is completion-independent.  Nothing in the construction knows which
completion it is in until the values $\zeta(i,\theta)$ or $\zeta_p(i,\theta)$ are
attached.

\subsection{What they do not share}

Against that shared skeleton, the constants produced are of wholly different
kinds.

\begin{proposition}[archimedean tameness]\label{prop:tame}
\IMPORTED.  For each of the twelve sporadic pairs whose characteristic roots have
distinct moduli, the archimedean Ap\'ery limit $\lim_nB(n)/A(n)$ exists and is a
rational multiple of $\zeta(w)$ or of $L(\chi_D,w)$ with $w\in\{2,3\}$
\cite{Zagier2009,ASvSZ2008,ChamberlandStraub,PaperA}.  In particular each spans a
one-dimensional $\QQ$-vector space, and the $\QQ$-span of all twelve is
contained in the five-dimensional space
$\QQ\zeta(2)+\QQ\zeta(3)+\QQ L(\chi_{-3},2)+\QQ L(\chi_{-4},2)+\QQ L(\chi_{-3},3)$.
(That the five are $\QQ$-linearly independent is of course not known; five is an
upper bound for the dimension, which is all we use.)
\end{proposition}

\begin{theorem}[the asymmetry]\label{thm:asym}
\PROVED, under Hypothesis~\ref{hyp:ind}, for the classical Ap\'ery pair.
Fix $p\ge5$ and $p\nmid a$.  The block decomposition \eqref{eq:Lamblocks} of the
single constant $\Lambda_a$ involves the tower Betas $\bet(j,ap^i-j)$ and
$\bet(j,ap^i)$ for every $i\ge0$ and every $j\le ap^i$ prime to $p$.  The group
$\Gamma_a$ they generate inside $1+p\Zp$ is free abelian of infinite rank.
\end{theorem}

\begin{proof}
Theorem~\ref{thm:Lam} exhibits the blocks.  By \eqref{eq:key} and
Hypothesis~\ref{hyp:ind}, $\Gamma_a$ is isomorphic to the image under $\delta$
of the free abelian group on those pairs; that image contains
$\mu_i:=\delta(ap^i,1)=[ap^i]-[1]-[ap^i-1]$ for every $i\ge0$.  The largest
symbol occurring in $\mu_i$ is $[ap^i]$, with coefficient $+1$, and the $ap^i$
are strictly increasing; so in any non-trivial finite $\ZZ$-combination
$\sum n_i\mu_i$ the symbol $[ap^{i_{\max}}]$ with $i_{\max}$ the largest index
with $n_{i}\ne0$ survives with coefficient $n_{i_{\max}}\ne0$.  Hence the $\mu_i$
are $\ZZ$-independent and the image has infinite rank
(Proposition~\ref{prop:image}, Corollary~\ref{cor:free}).
\end{proof}

The contrast is therefore not a matter of how hard the two sides are to compute.
The archimedean limit is a value of one period at one singular point; the
$p$-adic limit is a sum over the whole $p$-adic tower tree, and it collects the
generating function $G$ at \emph{every} level.  A one-dimensional space on one
side, an infinite-rank one on the other, with the rational skeleton identical.
Only the values attached differ, and $\zeta_p$ is attached infinitely often
where $\zeta$ is attached once.

\begin{remark}[the limitless trio]
Three of the fifteen families --- $B$, $(\delta)$, $(\eta)$ --- have
characteristic roots that are complex conjugates of equal modulus and hence
\emph{no} archimedean Ap\'ery limit at all, yet all three have $p$-adic tower
limits: $(\delta)$ outright, $B$ and $(\eta)$ after the $\chi(p)^{-s}$ twist
\cite[\S7]{PaperC}, the sole exception being $(\eta)$ at the ramified prime
$p=5$.  Proposition~\ref{prop:tame} is thus vacuous for them while
Theorem~\ref{thm:asym}'s $p$-adic side is not: on those three families the
asymmetry is total.  (The observation that the $p$-adic limit exists where the
archimedean one does not is due to \cite[\S8.2]{PaperA}; what is added here is
the algebra it lives in.)
\end{remark}

% =====================================================================
\section{The exclusion ledger}\label{sec:ledger}
% =====================================================================

Everything conditional above rests on Hypothesis~\ref{hyp:ind}.  We cannot prove
it; we can and do report what a search finds.

\subsection{Instrument and convention}

Reference values $\zeta_p(m)$, for all odd $m$ up to a cutoff a few beyond the
target precision (the $m$-th term of $\gh$ has valuation $\ge\varepsilon_m+m$ by
\eqref{eq:eps}), come from the Kubota--Leopoldt
implementation of \cite[App.~A]{PaperC}, built from a Volkenborn series and
validated on $980$ independent interpolation identities, on independence of the
auxiliary modulus, and against the limit characterisation of
\cite{LSZ2025,BV2025}.  The search instrument is an exact-\texttt{Fraction} LLL
with the residue column weighted by $p^{A}$, so that only exact relations modulo
$p^A$ survive.  For an $m$-term relation certified to $A$ digits the heuristic
noise floor is $p^{A/m}$: a lattice of that determinant contains vectors of that
size for accidental reasons.  LLL does not \emph{exclude} relations: it fails to
find them, with an approximation guarantee.  What we therefore report is: no
vector below the floor was returned, and LLL returns a vector within a factor
$2^{(n-1)/2}$ of the shortest, so a genuine relation substantially below the
floor would have been returned.  The decisive test is the second one:
\emph{a genuine relation's height does not move when $A$ increases}, whereas
noise tracks the floor.  All of this is evidence, not exclusion.

\begin{finding}\label{find:instrument}
\VER{$p\in\{5,7,11,13\}$; at the same precision as the searches below}  Adjoin
to the basis the number $y:=3\gh(1)-7\gh(2)$, computed from the same data.  The
finder recovers the planted relation $3\gh(1)-7\gh(2)-y=0$, of height $7$, at
every prime tested.
\end{finding}

\subsection{The search}

\begin{finding}[nothing below the floor]\label{find:ind}
\VER{$p\in\{5,7,11,13\}$; $x=1,\dots,8$; at $A=26$ and again at $A=38$ certified
$p$-adic digits}  No $\QQ$-linear relation among $\gh(1),\dots,\gh(N)$, and no
$\QQ$-affine one (i.e.\ with $1$ adjoined), was found below the noise floor, for
any $N\le8$.  Every best height sits at the floor and \emph{scales with the
precision}, which is the signature of noise.  Table~\ref{tab:ind} records the
two-precision comparison at $p=5$ and $p=7$; the other two primes behave
identically.
\end{finding}

\begin{table}[h]
\centering\small
\caption{$(\mathrm{IND}_p)$: best LLL height found against the noise floor
$p^{A/m}$, at two certified precisions.  A genuine relation would have a height
that does not move between the two columns.}
\label{tab:ind}
\begin{tabular}{@{}lcccc@{}}
\toprule
& \multicolumn{2}{c}{$p=5$} & \multicolumn{2}{c}{$p=7$}\\
\cmidrule(lr){2-3}\cmidrule(lr){4-5}
basis & $A=26$ & $A=38$ & $A=26$ & $A=38$\\
      & height / floor & height / floor & height / floor & height / floor\\
\midrule
$\gh(1),\gh(2)$          & 8.6e8 / 1.2e9 & 1.9e13 / 1.9e13 & 7.4e10 / 9.7e10 & 1.0e15 / 1.1e16\\
$\gh(1..3)$              & 3.1e5 / 1.1e6 & 3.8e8 / 7.1e8   & 8.1e6 / 2.1e7   & 2.2e10 / 5.1e10\\
$\gh(1..4)$              & 1.1e4 / 3.5e4 & 2.3e6 / 4.4e6   & 1.6e5 / 3.1e5   & 7.4e7 / 1.1e8\\
$\gh(1..6)$              & 4.5e2 / 1.1e3 & 1.5e4 / 2.7e4   & 2.7e3 / 4.6e3   & 8.9e4 / 2.3e5\\
$\gh(1..8)$              & 95 / 1.9e2    & 9.2e2 / 2.1e3   & 2.5e2 / 5.6e2   & 4.0e3 / 1.0e4\\
\midrule
$1,\gh(1)$               & 8.1e8 / 1.2e9 & 1.2e13 / 1.9e13 & 7.2e10 / 9.7e10 & 1.0e16 / 1.1e16\\
$1,\gh(1..3)$            & 7.9e3 / 3.5e4 & 2.2e6 / 4.4e6   & 1.1e5 / 3.1e5   & 3.3e7 / 1.1e8\\
$1,\gh(1..8)$            & 58 / 1.1e2    & 3.6e2 / 8.9e2   & 1.5e2 / 2.8e2   & 1.7e3 / 3.7e3\\
\bottomrule
\end{tabular}
\end{table}

\begin{remark}[what the ledger does and does not reach]
It reaches heights up to roughly $10^{13}$ at $p=5$ and $10^{15}$ at $p=7$ in the
two-term case, degrading to a few hundred for eight terms, at $x\le8$; no
relation was returned below the floor anywhere.  It says nothing about
relations of larger height, relations
involving $\gh$ at larger arguments, or --- crucially --- relations that are not
of the form $\widehat G(\mu)=0$ for $\mu$ of small support.  Under
Theorem~\ref{thm:rigid} the relations that matter for the Kazandzidis algebra are
exactly those with $\mu\in\im\delta=\ker m_1$, so the ledger is a search in the
right place; it is not a proof of anything.
\end{remark}

\subsection{What is imported}

Paper~C's exclusion theorem \cite[\S5]{PaperC} is the companion negative and we
do not repeat its tables: no $\QQ$-affine relation of height $\lesssim10^4$
between $\Lambda_a$ (or the Brown--Zudilin analogues) and $\zeta_p(3)$ or
$\zeta_p(5)$, at $p=5,7,11,13$, to $12$--$27$ certified digits; the same for
$\log_p2$, $\log_p3$, the unit root of $\eta(2z)^4\eta(4z)^4$, and for
algebraicity of degree $\le3$; and, across families, no pairwise $\QQ$-linear
relation among the $\Lambda_1$'s (at $p=7$, $12$--$21$ certified digits).

Theorem~\ref{thm:rigid} does \emph{not} apply to $\Lambda_a$, which is an
infinite assembly of blocks rather than a finite monomial in tower Betas, and we
claim no implication.  What the present paper adds to those exclusions is a
heuristic, and we label it as one: by Theorems~\ref{thm:closed} and
\ref{thm:Lam} the ingredients of $\Lambda_a$ are exponentials of infinite
$\QQ$-combinations of \emph{all} the odd $\zeta_p(m)$ with $p$-power-weighted
coefficients, so a finite low-height relation between $\Lambda_a$ and a single
$\zeta_p(m)$ would be a coincidence of the same species as the ones
Hypothesis~\ref{hyp:ind} forbids --- but it is not one of them, and nothing here
proves it cannot happen.

% =====================================================================
\section{The programme}\label{sec:programme}
% =====================================================================

The $\Lambda$-algebra is a class of constants with closed forms and no theory.
Here are the questions we can now state precisely, in the order we would attack
them.

\begin{problem}[the presentation]\label{prob:present}
Prove Problem~\ref{prob:gen}: $\ker\delta$ is generated by the symmetry and
trinomial-revision relations.  This is a finite-type statement about the free
abelian group $\ker m_1$ and is independent of $p$; Finding~\ref{find:lattice}
confirms it to $N=11$.  With it, Theorem~\ref{thm:rigid} becomes a complete
determination of the multiplicative relations among the tower Betas, conditional
only on $(\mathrm{IND}_p)$.  (For the caveat about $\El$ itself, see
Remark~\ref{rem:notclaimed}.)
\end{problem}

\begin{problem}[$(\mathrm{IND}_p)$ in weakened form]\label{prob:indweak}
Hypothesis~\ref{hyp:ind} in full is at least as hard as the irrationality of a
single $\zeta_p(2k+1)$ for $p\ge5$.  Two weakenings look accessible:
\emph{(a)} prove $\gh(1)\notin\QQ$ for some $p$ --- equivalently, that
$\sum_{m\ \mathrm{odd}\ge3}\zeta_p(m)p^m/(m(1-p^m))$ is irrational;
\emph{(b)} prove that $\gh(1),\dots,\gh(N)$ span a $\QQ$-space of dimension
$\ge2$ for some explicit $N$ and $p$.  Both are $p$-adic Ap\'ery-type questions
about a \emph{single} very rapidly convergent series, and (b) is the natural
target for a $p$-adic Ball--Rivoal argument \cite{LaiSprang2023,FSZ2018}.
\end{problem}

\begin{problem}[the $\Lambda$-module across bases]\label{prob:bases}
Determine the $\QQ$-vector space spanned by $\{\Lambda_a\}_{p\nmid a}$ inside
$\Lalg$.  Paper~C excludes low-height relations; Theorem~\ref{thm:Lam} shows the
$a$-dependence enters through $h_p(a)$, through $\Atil_a$ and through the block
index set $\{j\le ap^i\}$.  The tractable sub-question: is
$\Lambda_a-h_p(a)$ determined by $\Lambda_1-h_p(1)$ and finitely many
$\Gp$-values?  The block structure says no, and a proof would be the first
unconditional non-relation in this subject.
\end{problem}

\begin{problem}[the other fourteen families]\label{prob:families}
Theorems~\ref{thm:Atil} and \ref{thm:Lam} are proved for the classical Ap\'ery
$\zeta(3)$ pair, using Ap\'ery's explicit harmonic formula for $b_n$.  Every
sporadic $B$ has a harmonic-monomial closed form of weight $w$
\cite[\S6]{PaperA}; carry the same two estimates through it and one gets
$\Lambda_a\in\Lalg$ for all fifteen families, with $\El$ enlarged by the
character values.  What is the correct twisted analogue of $h^{(w)}_p$, and does
$\chi$ enter the blocks or only the normalisation?
\end{problem}

\begin{problem}[a weight filtration]\label{prob:weight}
Is there a filtration on $\Lalg$ under which $h^{(w)}_p(a)$ has weight $w$,
$\bet(B,D)$ has weight $0$, and $\Lambda_a$ has weight $w$?  The obstruction is
that $\bet$ is an exponential of a series mixing all odd weights; a filtration
would have to be by the \emph{order of vanishing} of $\gh$ at $x=0$, which is
$3$ --- the same for every tower Beta, so the filtration would be trivial.  A graded object, if one exists, is the right home for a
motivic statement.
\end{problem}

\begin{problem}[what a period statement would mean]\label{prob:period}
The archimedean Ap\'ery limits are periods.  Their $p$-adic counterparts are
values of $\prod_{k\ge1}G(xp^k)$ with $G=\Gp e^{-\Gp'(0)x}$, the same generating
function whose Taylor coefficients are the Beukers--Vlasenko Frobenius structure
constants \cite{BV2025}; archimedeanly the analogous generating function is the
one of Bloch--Vlasenko \cite{BlochVlasenko}.  Is there a category in which
$\prod_{k\ge1}G(xp^k)$ is the period of an object, and in which the divisor map
$\delta$ of Definition~\ref{def:delta} is the shadow of a coproduct?  The
relations \eqref{eq:relS}, \eqref{eq:relT} are exactly those of a
$\ZZ$-valued ``divisor of a binomial'', which is suggestive.
\end{problem}

\begin{problem}[falsifiable next computations]\label{prob:next}
Three, in increasing cost.
\begin{enumerate}[label=(\arabic*),leftmargin=2em]
\item Extend Finding~\ref{find:lattice} to $N=20$; a failure there would refute
  Problem~\ref{prob:present} outright.  Cost: exact Smith normal forms of
  matrices of size $\binom{20}{2}=190$.
\item Extend Finding~\ref{find:ind} to $x\le16$ and $A=60$ certified digits.
  $\gh$ costs only $\zeta_p(m)$ for $m\le A+6$, so this is cheap; the
  discriminating quantity is whether any height fails to scale with $A$.
\item Compute $\Lambda_a$ for the twisted families to $18$ digits and test
  $\Lambda^{(\mathrm{fam}_1)}_1/\Lambda^{(\mathrm{fam}_2)}_1$ against
  $\El$-monomials rather than against $\QQ$ --- Paper~C tested the latter only,
  and Theorem~\ref{thm:rigid} says $\El$ is the right group to quotient by.
\end{enumerate}
\end{problem}

% =====================================================================
\appendix
\section{Reproduction}\label{app:repro}
% =====================================================================

All computations use exact integer or rational arithmetic; no floating point
occurs at any stage.  The programs accompanying this paper are:

\begin{itemize}[leftmargin=1.5em]
\item \texttt{lam.py} --- the $p$-adic number class with an honest precision
  bound carried through every operation; $\gh$, $\bet$ and $E_p$ from the
  Kubota--Leopoldt values $\zeta_p(m)$; the direct Kazandzidis limits
  $c_p(A,B)$ from exact binomial residues by Legendre/Morita splitting (one
  linear pass per prime, recording only the needed factorial-unit indices); the
  verification of \eqref{eq:closed} (Finding~\ref{find:closedver}); the three
  relation sweeps (Finding~\ref{find:rels}); the $(\mathrm{IND}_p)$ hunt with the
  planted-relation validation (Findings~\ref{find:instrument},
  \ref{find:ind}); and the exact integer linear algebra for
  Finding~\ref{find:lattice}, including the assertion that every proposed
  generator really lies in $\ker\delta$; and the tower-harmonic constants
  $h^{(w)}_p(a)$, with the distribution relation and the relation hunt of
  Finding~\ref{find:hp} (the index $i$ ranging over all integers with
  $ap^i\ge1$, cf.\ Remark~\ref{rem:index}).
\item \texttt{atil.py} --- the block-by-block check of Theorem~\ref{thm:Atil}
  for $a=1,2$ (Finding~\ref{find:blocks}), with $a_{ap^s}$ built independently as
  an exact integer from the defining double sum.
\end{itemize}

The $\zeta_p(m)$ implementation, its validation record and the exact LLL are
those of \cite[App.~A]{PaperC} and are re-used unchanged; the certified
precision of every reported digit is carried by the number class rather than
asserted.

\subsection*{Acknowledgements}

[TO BE SUPPLIED]

% =====================================================================
\begin{thebibliography}{99}
% =====================================================================

% --- Entries marked (BV) were resolved against a fetched source in the
%     verification memo work/BIBLIO_VERIFY.md of the companion series; the
%     remainder were fetched for this paper.

\bibitem{PaperA}
[AUTHOR],
\emph{$\chi$-twisted Lucas congruences for the second solutions of Ap\'ery-like
recurrences}.
Companion paper (Paper~A of this series); in preparation.

\bibitem{PaperB}
[AUTHOR],
\emph{Frobenius constants, Stokes ratios, and the middle-root obstruction for
cellular approximations to $\zeta(5)$}.
Companion paper (Paper~B of this series); in preparation.

\bibitem{PaperC}
[AUTHOR],
\emph{$p$-adic Ap\'ery limits: rate laws, valuation strata, and closed forms via
the $p$-adic Gamma function}.
Companion paper (Paper~C of this series); in preparation.

\bibitem{PaperD}
[AUTHOR],
\emph{Rhin--Viola groups are $p$-adic: the transfer principle, the margin map,
and the optimality of the known $p$-adic irrationality measures}.
Companion paper (Paper~D of this series); in preparation.

\bibitem{Apery1979}
R.~Ap\'ery,
\emph{Irrationalit\'e de $\zeta(2)$ et $\zeta(3)$},
in: Journ\'ees Arithm\'etiques (Luminy, 1978),
Ast\'erisque \textbf{61} (1979), 11--13.

\bibitem{ASvSZ2008}
G.~Almkvist, D.~van Straten and W.~Zudilin,
\emph{Ap\'ery limits of differential equations of order $4$ and $5$},
in: Modular Forms and String Duality, Fields Inst.\ Commun.\ \textbf{54},
Amer.\ Math.\ Soc., Providence, RI, 2008, pp.~105--123.

\bibitem{BlochVlasenko}
S.~Bloch and M.~Vlasenko,
\emph{Gamma functions, monodromy and Frobenius constants},
Commun.\ Number Theory Phys.\ \textbf{15} (2021), no.~1, 91--147;
arXiv:1908.07501.

\bibitem{BV2025}
F.~Beukers and M.~Vlasenko,
\emph{Frobenius structure and $p$-adic zeta values},
Adv.\ Math.\ \textbf{480} (2025), Paper No.~110512; arXiv:2302.09603.
(Equation~(2) there is $\log\Gp(x)=\Gp'(0)x-\sum_{m\ge2}\zeta_p(m)x^m/m$, read
verbatim from the fetched arXiv text.  The evaluation of the Frobenius structure
constants as Taylor coefficients of $\Gp(x)e^{-\Gp'(0)x}$ is Theorem~1.4 in that
text and is cited as Theorem~1.5 in \cite{PaperC}; we therefore cite it without
an internal number.)

\bibitem{Calegari2005}
F.~Calegari,
\emph{Irrationality of certain $p$-adic periods for small $p$},
Int.\ Math.\ Res.\ Not.\ (2005), no.~20, 1235--1249.

\bibitem{ChamberlandStraub}
M.~Chamberland and A.~Straub,
\emph{Ap\'ery limits: experiments and proofs},
Amer.\ Math.\ Monthly \textbf{128} (2021), no.~9, 811--824;
arXiv:2011.03400.

\bibitem{Cooper2012}
S.~Cooper,
\emph{Sporadic sequences, modular forms and new series for $1/\pi$},
Ramanujan J.\ \textbf{29} (2012), no.~1--3, 163--183.
(The three additional sporadic solutions $s_7$, $s_{10}$, $s_{18}$ are
attributed to this paper by \cite[\S2]{MalikStraub}.)

\bibitem{Fischler2003}
S.~Fischler,
\emph{Irrationalit\'e de valeurs de z\^eta (d'apr\`es Ap\'ery, Rivoal, \dots)},
S\'eminaire Bourbaki exp.\ 910, Ast\'erisque \textbf{294} (2004), 27--62;
arXiv:math/0303066.

\bibitem{FSZ2018}
S.~Fischler, J.~Sprang and W.~Zudilin,
\emph{Many odd zeta values are irrational},
Compositio Math.\ \textbf{155} (2019), no.~5, 938--952; arXiv:1803.08905.

\bibitem{Granville1997}
A.~Granville,
\emph{Arithmetic properties of binomial coefficients.\ I.\ Binomial coefficients
modulo prime powers},
in: Organic Mathematics (Burnaby, BC, 1995), CMS Conf.\ Proc.\ \textbf{20},
Amer.\ Math.\ Soc., Providence, RI, 1997, pp.~253--276.

\bibitem{GrossKoblitz1979}
B.~H.~Gross and N.~Koblitz,
\emph{Gauss sums and the $p$-adic $\Gamma$-function},
Ann.\ of Math.\ (2) \textbf{109} (1979), no.~3, 569--581.

\bibitem{JacobsthalKazandzidis}
V.~Brun, J.~O.~Stubban, J.~E.~Fjeldstad, R.~Tambs Lyche, K.~E.~Aubert,
W.~Ljunggren and E.~Jacobsthal,
\emph{On the divisibility of the difference between two binomial coefficients},
Den 11te Skandinaviske Matematikerkongress (Trondheim, 1949),
Johan Grundt Tanums Forlag, Oslo, 1952, pp.~42--54;
G.~S.~Kazandzidis,
\emph{Congruences on the binomial coefficients},
Bull.\ Soc.\ Math.\ Gr\`ece (N.S.) \textbf{9} (fasc.~1) (1968), 1--12.

\bibitem{LaiSprang2023}
L.~Lai and J.~Sprang,
\emph{Many $p$-adic odd zeta values are irrational},
arXiv:2306.10393.

\bibitem{LSZ2025}
L.~Lai, J.~Sprang and W.~Zudilin,
\emph{A note on the irrationality of $\zeta_2(5)$},
arXiv:2505.05005.
(Definition~8 there fixes the normalisation $\zeta_p(s)=L_p(s,\omega^{1-s})$
used here.)

\bibitem{MalikStraub}
A.~Malik and A.~Straub,
\emph{Divisibility properties of sporadic Ap\'ery-like numbers},
arXiv:1508.00297.

\bibitem{Mestrovic}
R.~Me\v{s}trovi\'c,
\emph{Wolstenholme's theorem: its generalizations and extensions in the last
hundred and fifty years (1862--2012)},
arXiv:1111.3057.
(A survey; \S3 states the Jacobsthal--Kazandzidis congruence
$\bin{np}{mp}\equiv\bin nm\pmod{p^{t}}$, $t=\vp\bigl(p^3nm(n-m)\bigr)$, and its
iterate, which is the classical source of the convergence in \eqref{eq:cp}.)

\bibitem{Morita1975}
Y.~Morita,
\emph{A $p$-adic analogue of the $\Gamma$-function},
J.\ Fac.\ Sci.\ Univ.\ Tokyo Sect.\ IA Math.\ \textbf{22} (1975), no.~2,
255--266.
(The product form in \eqref{eq:morita-int} is the iterate of Morita's relation
$\Gp(n+1)=(-1)^{n+1}n!/(p^{\floor{n/p}}\floor{n/p}!)$.)

\bibitem{Robert2000}
A.~M.~Robert,
\emph{A Course in $p$-adic Analysis},
Graduate Texts in Mathematics \textbf{198}, Springer-Verlag, New York, 2000,
xv$+$437 pp.
(A standard reference for Morita's $\Gp$, its Lipschitz continuity and the
relation $\Gp(x+1)=-\Gp(x)$ on $p\Zp$; these are also in \cite{Morita1975}.)

\bibitem{Straub2023}
A.~Straub,
\emph{Gessel--Lucas congruences for sporadic sequences},
Monatsh.\ Math.\ (2023); arXiv:2301.12248.

\bibitem{Supercong}
F.~Beukers,
\emph{Some congruences for the Ap\'ery numbers},
J.\ Number Theory \textbf{21} (1985), no.~2, 141--155;
M.~J.~Coster,
\emph{Supercongruences},
Ph.D.\ thesis, Universiteit Leiden, 1988;
T.~Ishikawa,
\emph{Super congruence for the Ap\'ery numbers},
Nagoya Math.\ J.\ \textbf{118} (1990), 195--202.

\bibitem{Zagier2009}
D.~Zagier,
\emph{Integral solutions of Ap\'ery-like recurrence equations},
in: J.~Harnad and P.~Winternitz (eds.), Groups and Symmetries:
From Neolithic Scots to John McKay,
CRM Proc.\ Lecture Notes \textbf{47}, Amer.\ Math.\ Soc., 2009, pp.~349--366.

\end{thebibliography}

\end{document}
